\begin{figure}[ht]
\centering
\begin{tikzpicture}[x=0.55cm, y=0.45cm, font=\footnotesize]

% x-axis: log10 N (number of samples / evaluations), 1 to 8
% y-axis: log10 |error|, 0 down to -16
\draw[->, thick] (0,0) -- (9,0) node[right] {$\log_{10} N$};
\draw[->, thick] (0,-16) -- (0,1) node[above] {$\log_{10}\abs{\text{error}}$};

\foreach \k in {1,2,3,4,5,6,7,8} {
  \draw (\k,0) -- (\k,0.3);
  \node[anchor=south, font=\scriptsize] at (\k,0.35) {\k};
}
\foreach \d in {0,-4,-8,-12,-16} {
  \draw (0,\d) -- (-0.2,\d);
  \node[anchor=east, font=\scriptsize] at (-0.2,\d) {$10^{\d}$};
}

% Monte Carlo: error ~ C / sqrt(N), slope -1/2 on log-log, with jitter.
\draw[thick, engred] plot[smooth] coordinates {
  (1,-0.6) (2,-0.9) (3,-1.7) (4,-2.1) (5,-2.4) (6,-3.2) (7,-3.4) (8,-4.1)
};
\node[engred, anchor=west, font=\scriptsize] at (6.1,-2.7) {Monte Carlo $\sim N^{-1/2}$};

% Trapezoidal (1D): slope -2 in N (h ~ 1/N)
\draw[thick, engblue] (1,-1) -- (7,-13);
\node[engblue, anchor=west, font=\scriptsize] at (4.4,-8.8) {trapezoidal $\sim N^{-2}$};

% Gauss (1D smooth): exponential, reaches floor by N ~ 20 (log10 ~ 1.3)
\draw[thick, engamber] plot coordinates {
  (1,-4) (1.3,-16)
};
\draw[thick, engamber] (1.3,-16) -- (3,-16);
\node[engamber!80!black, anchor=west, font=\scriptsize] at (1.5,-15.1) {Gauss (smooth 1D)};

\draw[thick, dashed, enggray] (0,-16) -- (9,-16);

% dimension arrow: MC slope is dimension-independent
\node[align=center, font=\scriptsize, anchor=north] at (4.5,-15.4)
  {Monte Carlo slope is $-1/2$ in \emph{any} dimension};

\end{tikzpicture}
\caption[Monte Carlo versus deterministic quadrature]{Integration error
against the number of function evaluations $N$, on log-log axes. In one
dimension the deterministic rules dominate: the trapezoidal rule falls at
slope $-2$ and a Gauss rule converges geometrically to the round-off
floor within a few dozen evaluations. Monte Carlo estimation, averaging
the integrand at $N$ random points, falls only at slope $-1/2$ set by the
$\sigma/\sqrt{N}$ standard error of the sample mean, and its individual
realisations scatter about that trend. The decisive property is the one
the plot cannot show in one dimension: the Monte Carlo slope is $-1/2$
regardless of the dimension, while a product deterministic rule needs
$N = m^{d}$ points for $m$ per axis in $d$ dimensions and becomes
useless past $d \approx 6$. Monte Carlo is the working choice for
high-dimensional integrals precisely because its slow convergence does
not degrade with dimension. Source: authors' analysis.}
\label{fig:vol02ch16:monte-carlo-convergence}
\end{figure}
